\newif\ifglobal

%%%%%%%%%%%%%%%%%%%%%%%
%%% Compile to view %%%
%%%%%%%%%%%%%%%%%%%%%%%

%%%%%%%%%%%%%%%%%%%%%%%%%%%%%%%%%%%%%%%%%%%%%%%%%%%
%%% Readme:                                     %%%
%%% https://www.overleaf.com/read/bwdjvfjksvjy/ %%%
%%%%%%%%%%%%%%%%%%%%%%%%%%%%%%%%%%%%%%%%%%%%%%%%%%%

% >>> M?: marks a module setting number or important notice

\def \modulefiles {./20-RTCTIMER}

% >>> M1: This module is in global project? <<<
% 'YES' - uncomment; 'NO' - comment out
\globaltrue

\ifglobal
    \documentclass[main\_\_EN.tex]{subfiles}

\else
    \documentclass[openany, 10pt]{book}

    % >>> M2: In ./chip-spec-settings/preamble-trm-module, also find
    %         settings for visibility of todo notes, watermark <<<
    \usepackage{preamble-trm-module}

    % >>> M3: Temporary labels in English,
    %         you can add your own labels there <<<
    \myexternaldocument{temp-labels-en}

\fi %\ifglobal


\setcounter{secnumdepth}{4}

% >>> M4: Set {./relative-path} to external files for this module <<<
% To add an external file, use:
% (Replace '---' with actual filename)
% '\input{\modulefiles/tables/---.tex}' for tables
% '\includegraphics{\modulefiles/figures/---}' for figures
% '\subfile{\modulefiles/---}' for register files


% >>> M5: If this module needs any updates in
% './chip-spec-settings/preamble-trm-module.sty'
% (include additional package, etc.), contact your module PM
% or create the file './chip-spec-settings/changelog.tex' make the
% necessary updates yourself and document them in this file <<<

\raggedright


\begin{document}


% Sets language into which pre-defined bits of text are translated
\selectlanguage{english}

% Do not delete this block
\ifglobal
% Add the following front matter if
% this module is outside of global project
\else
    \InputIfFileExists{readme}{}{}
    {\let\clearpage\relax \listoftodos}
    {\let\clearpage\relax \listofdonetodos}
    \tableofcontents
    %\listoftables
    %\listoffigures
\fi


%%%%%%%%%%%%%%%%%%%%%%%%%
%%% TRM Module Begins %%%
%%%%%%%%%%%%%%%%%%%%%%%%%

\hypertarget{rtctimer}{}
\chapter{RTC Timer}
\label{mod:rtctimer}

%%%%%%%%%%%%%%%%%%%%
%%% Introduction %%%
%%%%%%%%%%%%%%%%%%%%

\section{Introduction}

RTC Timer is an important module for implementing low power management of \chipname{}. Based on a 48-bit readable counter, RTC Timer is mainly used as a system timer in low power mode when the timer peripheral in the HP system is unavailable. It also allows for configuring timer interrupts and logging the time when specific events happen in the system. The supported events are described in the \ref{sec:rtctimer-func-descr} \nameref{sec:rtctimer-func-descr}.

\section{Feature List} \label{sec:rtctimer-features}

RTC Timer has the following features:


\begin{itemize}
    \item 48-bit counter
    \item Time logging when one of the following events happens:
    \begin{itemize}
        \item HP system reset
        \item CPU enters stall state
        \item CPU exits stall state
        \item Crystal powers up
        \item Crystal powers down
    \end{itemize}
    \item Time logging through register configuration
    \item Occurrence time cached of the most recent two specific events
    \item Generation of interrupts at target times, which are configurable. It is also possible to configure two target times simultaneously.
    \item Uninterrupted operation during any reset or sleep mode, except for power-on reset of LP system.
\end{itemize}

%%%%%%%%%%%%%%%%%%%%%%%%%%%%%%
%%% Functional Description %%%
%%%%%%%%%%%%%%%%%%%%%%%%%%%%%%

\section{Functional Description}
\label{sec:rtctimer-func-descr}

\begin{itemize}
    \item 48-bit counter
    \begin{itemize}
        \item The implementation of the RTC Timer is based on a 48-bit counter, driven by the RC\_SLOW\_CLK in the Always\-On power domain. It uses continuous loop counting (except during LP system reset), and an overflow interrupt is generated when the 48-bit counter overflows.
    \end{itemize}


    \item Log the occurrence time of specific events
    \begin{itemize}
        \item RTC Timer supports logging the occurrence time of three types of events:
        \begin{itemize}
            \item HP system reset
            \item CPU enters or exits stall state
            \item Crystal powers up or down
        \end{itemize}
        \item The above three types of events have their own independent configuration registers. When the configuration registers are enabled, if the corresponding events occur, the hardware pulse generated will trigger the RTC Timer to log the current time. The corresponding configuration registers are shown in the table below.

        \begin{longtable}[c]{|p{6cm}|p{9cm}|}
        \caption{Configure RTC Timer to Log the Occurrence Time of Specific Events} \label{tab:configure-rtc-timer-log-time-of-specific-events}
        \\\hline
        \rowcolor{lightgray}
        \textbf{Configuration Register} & \textbf{Description} \\
        \hline
        \hyperref[regdesc:RTCTIMERUPDATEREG]{RTC\_TIMER\_MAIN\_TIMER\_SYS\_RST} &  Enable the RTC Timer to log the time of system reset. \\
        \hline
        \hyperref[regdesc:RTCTIMERUPDATEREG]{RTC\_TIMER\_MAIN\_TIMER\_SYS\_STALL} & Enable the RTC Timer to log the time when the CPU enters or exits the stall state.\\
        \hline
        \hyperref[regdesc:RTCTIMERUPDATEREG]{RTC\_TIMER\_MAIN\_TIMER\_XTAL\_OFF} &  Enable the RTC Timer to log the time when the crystal powers up or down.\\
        \hline
        \end{longtable}
    \end{itemize}


    \item Log the current time through software configuration.
    \begin{itemize}
        \item In addition to the above three specific events that can trigger the RTC Timer to log the time, the RTC Timer can also record the current time through the configuration register \hyperref[regdesc:RTCTIMERUPDATEREG]{RTC\_TIMER\_MAIN\_TIMER\_UPDATE}.
        \item A pulse signal will be generated by configuring the \hyperref[regdesc:RTCTIMERUPDATEREG]{RTC\_TIMER\_MAIN\_TIMER\_UPDATE}, which will trigger the RTC Timer to record the current time.
    \end{itemize}

    \item Generate counting interrupts at target times
    \begin{itemize}
        \item RTC Timer supports configuring two target times simultaneously, referred to as target time 0 and target time 1. When the timer reaches target time 0, it will trigger the RTC\_TIMER\_CMP0\_INT interrupt. Similarly, when the timer reaches Target Moment 1, it will trigger the RTC\_TIMER\_CMP1\_INT interrupt.

        \item Configure the target times through the registers below:

        \begin{longtable}[c]{|p{6cm}|p{9cm}|}
        \caption{Target times Configuration} \label{tab:rtc-timer-target-time-configuration}
        \\\hline
        \rowcolor{lightgray}
        \textbf{Configuration Register} & \textbf{Description} \\
        \hline
        \multirow{1}{*}{\hyperref[regdesc:RTCTIMERTAR0HIGHREG]{RTC\_TIMER\_MAIN\_TIMER\_TAR\_EN\regindex{n}}} & \regindex{n} = 0 or 1, enable target time configuration. \\
        \hline
        \multirow{1}{*}{\hyperref[regdesc:RTCTIMERTAR0LOWREG]{RTC\_TIMER\_MAIN\_TIMER\_TAR\_LOW\regindex{n}}} & \regindex{n} = 0 or 1, configure the low 32 bits of the target time.\\
        \hline
        \multirow{1}{*}{\hyperref[regdesc:RTCTIMERTAR0HIGHREG]{RTC\_TIMER\_MAIN\_TIMER\_TAR\_HIGH\regindex{n}}} & \regindex{n} = 0 or 1, configure the high 16 bits of the target time.\\
        \hline
        \end{longtable}
    \end{itemize}


    \item Read the time cached in the RTC Timer
    \begin{itemize}
        \item There are two sets of registers used to cache the occurrence time of specific events, namely register group 0 and register group 1.
        \item Register group 0 includes register \hyperref[regdesc:RTCTIMERMAINBUF0LOWREG]{RTC\_TIMER\_MAIN\_BUF0\_LOW\_REG} and \hyperref[regdesc:RTCTIMERMAINBUF1HIGHREG]{RTC\_TIMER\_MAIN\_BUF1\_HIGH\_REG} which are used to cache the count value of the RTC Timer under the current trigger.
        \item Register group 1 includes register \hyperref[regdesc:RTCTIMERMAINBUF1LOWREG]{RTC\_TIMER\_MAIN\_BUF1\_LOW\_REG} and \hyperref[regdesc:RTCTIMERMAINBUF1HIGHREG]{RTC\_TIMER\_MAIN\_BUF1\_HIGH\_REG}, which are used to cache the count value of the RTC Timer from the last trigger.
        \item On a new trigger, the record from the previous trigger will be moved from register group 0 to register group 1 (and the previous record in register group 1 will be overwritten), and the record of this trigger will be stored in register group 0. Therefore, only the last two triggers can be logged at any time.

    \end{itemize}
\end{itemize}


%%%%%%%%%%%%%%%%%%%%%%%%%%%%%%%%%
%%% Event Task Matrix Feature %%%
%%%%%%%%%%%%%%%%%%%%%%%%%%%%%%%%%

\section{Event Task Matrix Feature}
\label{sec:rtctimer-etm-func}


The RTC Timer on \chipname{} supports the Event Task Matrix (ETM) function, which allows RTC Timer's ETM tasks to be triggered by any peripherals' ETM events, or RTC Timer's ETM events to trigger any peripherals' ETM tasks. This section introduces the ETM tasks and events related to RTC Timer. For more information, please refer to Chapter \ref{mod:evntaskmatrix} \nameref{mod:evntaskmatrix}.

RTC Timer doesn't receive any ETM tasks:

RTC Timer can generate the following ETM events:

\begin{itemize}
    \item RTC\_EVT\_CMP: Indicate that the RTC Timer reaches the target time 0 configured by \hyperref[regdesc:RTCTIMERTAR0LOWREG]{RTC\_TIMER\_MAIN\_TIMER\_TAR\_LOW0} and \hyperref[regdesc:RTCTIMERTAR0HIGHREG]{RTC\_TIMER\_MAIN\_TIMER\_TAR\_HIGH0}.
    \item RTC\_EVT\_OVF: Indicate the 48-bit counter overflow.
    \item RTC\_EVT\_TICK: Indicates the event that occurs when the RTC timer increments by 1.
\end{itemize}


%%%%%%%%%%%%%%%%%%%%%
%%% Interrupts %%%
%%%%%%%%%%%%%%%%%%%%%

\section{Interrupts}

\chipname{}'s RTC Timer can generate the following interrupt signal(s) that will be sent to the \hyperref[mod:intmtrx]{Interrupt Matrix}.

    \begin{itemize}
        \item LP\_TIMER\_REG\_0\_INTR
        \item LP\_TIMER\_REG\_1\_INTR
    \end{itemize}

There are several internal interrupt sources from RTC Timer that can generate the above interrupt signal(s). The interrupt sources from RTC Timer are listed with their trigger conditions and the resulted interrupt signal(s) in Table \ref{tab:rtctimer-int-sources}.


\begin{longtable}{ | L{5cm} | L{6.5cm} | L{4cm} | }
\caption{RTC Timer's Internal Interrupt Sources} \label{tab:rtctimer-int-sources}
\\ \hline
\rowcolor{lightgray}
\textbf{Internal Interrupt Source} & \textbf{Trigger Condition}  &  \textbf{Interrupt Signal}\\ \hline
\label{int:RTCTIMERCMP0INT}RTC\_TIMER\_CMP0\_INT &  RTC Timer reaches the target time 0 & LP\_TIMER\_REG\_0\_INTR  \\ \hline
\label{int:RTCTIMERCMP1INT}RTC\_TIMER\_CMP1\_INT &  RTC Timer reaches the target time 1 & LP\_TIMER\_REG\_1\_INTR \\ \hline
\label{int:RTCTIMEROVERFLOWINT}RTC\_TIMER\_OVERFLOW\_INT &  48-bit counter overflow & LP\_TIMER\_REG\_0\_INTR  LP\_TIMER\_REG\_1\_INTR\\ \hline
\end{longtable}

\begin{tiplisting}
For definitions of \textit{interrupt}, \textit{interrupt signal}, \textit{interrupt source}, and their correlations, please refer to Chapter \ref{mod:intmtrx} \textit{\nameref{mod:intmtrx}} > Section \ref{sec:intmtx-int-terminology} \textit{\nameref{sec:intmtx-int-terminology}}.
\end{tiplisting}

Each interrupt source can be configured by a common set of registers that are described in Section \hyperref[interrupt-config-registers]{\textit{Interrupt Configuration Registers}}. The specific registers can be found in Section \ref{sec:rtctimer-reg-summ} \textit{\nameref{sec:rtctimer-reg-summ}}.


% >>> If your module has no registers, delete sample text
%       for Sections Register Summary and Registers below <<<

%%%%%%%%%%%%%%%%%%%%%%%%
%%% Register Summary %%%
%%%%%%%%%%%%%%%%%%%%%%%%

% >>> Update the sample text below and insert
%     the info on register summary into the subfile <<<

\newpage
\hypertarget{rtctimer-reg-summ}{}
\section{Register Summary}
\label{sec:rtctimer-reg-summ}

The addresses in this section are relative to RTC Timer base address provided in Table \ref{tab:sysmem-base-address} in Chapter \ref{mod:sysmem} \textit{\nameref{mod:sysmem}}.

The abbreviations given in Column \textbf{Access} are explained in Section \hyperref[glossary-access-types]{\textit{Access Types for Registers}}.

\subfile{\modulefiles/20-rtctimer-reg-summ__EN}


%%%%%%%%%%%%%%%%%
%%% Registers %%%
%%%%%%%%%%%%%%%%%

% >>> Update the sample text below and insert
%      the info on registers into the subfile <<<

\newpage
\section{Registers}

The addresses in this section are relative to RTC Timer base address provided in Table \ref{tab:sysmem-base-address} in Chapter \ref{mod:sysmem} \textit{\nameref{mod:sysmem}}.

For how to program reserved fields, please refer to Section \hyperref[programming-reserved-field]{Programming Reserved Register Field}.

\subfile{\modulefiles/20-rtctimer-reg__EN}


\end{document}
